\section{Sampling the distribution with a \texorpdfstring{$\BPAC^0[1]$}{BPAC0[1]} circuit}\label{appendix:sample-bpac}

Let $p$ be the probability distribution computed by the distribution generator algorithm of a probabilistic, constant-depth, polynomial-embedding-dimension transformer.
In this section we show that sampling $p$ can be done with a constant-depth and polynomial-size random circuit using only a constant number of random gates.

\medskip

We can sample from the distribution $p$ by drawing a uniform random variable on the interval $[0,1)$.
We refer to the value of this variable as the seed.
This is because we can partition $[0,1)$ into disjoint intervals, one per token, each of length equal to the token's probability.
Let $p$ be the probability computed by the distribution generator. 
Let $\Sigma = \{\sigma_1, \dots, \sigma_{|\Sigma|}\}$.
Then the ranges are defined as $\text{range}(\sigma_i) = [\sum_{j=1}^{i-1}p(\sigma_j), \sum_{j=1}^{i}p(\sigma_j))$.

For example, if $\Sigma = \{t_1, t_2, t_3\}$ and $p$ is such that $p(t_1) = 0.33$, $p(t_2) = 0.22$, $p(t_3) = 0.45$, then the ranges would be $[0, 0.33)$ for $t_1$, $[0.33, 0.55)$ for $t_2$, and $[0.55, 1)$ for $t_3$.

Observe that although the numbers $p(\sigma_j)$ are representable in $\Floating$, the summations are not performed with iterative rounding.
They are computed with the maximum precision needed to avoid rounding error; this precision can be bounded as a function of $|\Sigma|$, $e$ and $s$ alone, independently of the values of $p(\sigma_i)$.

Since the precision for the numbers is fixed and the number of terms of the summations are bounded by the size of the alphabet, the ranges can be computed by circuits in $\AC^0$.

\medskip

Let us analyze the seed.
It is just a random number in the range $[0,1)$.
We can interpret it as a binary string.
How many bits of the seed are needed to determine which range it belongs to?
It suffices to examine as many bits of the seed as the logarithm of the smallest probability.

Since $p$ comes from the distribution generator, for all tokens, $p(\sigma)$ has to be representable in $\Floating$.
Then, the smallest non-empty range has size at least $\epsilon$, where $\epsilon$ is the smallest positive number representable in $\Floating$.
Let $r$ be the number of bits needed for representing $\epsilon$ in binary.
Therefore, knowing the first $r$ bits of the binary representation of the seed is enough to sample $p$.
Thus, we only need $r$ random bits for sampling $p$.

Note that $r$ only depends on the sizes ($e$ and $s$) of the significand and exponent of $\Floating$.
Thus, $r$ is constant for constant-precision transformers.

Observe that once the ranges are computed, they act as a lookup table with one entry for each number in $[0,1)$ whose binary representation has $r$ bits.
Since $r$ is a constant, the ranges can be interpreted as a constant-size lookup table.


Therefore, since the ranges can be computed by $\AC^0$ circuits and the lookup table has constant size, sampling from $p$ can be performed by a circuit in $\BPAC^0[1]$.